% Section V-B: Characterize (new framework; replaces Six_Bottleneck_States_Suffice.tex and
% Balance_Is_Not_a_Monotone_Proxy_for_Quality.tex, which are kept unchanged).
% Figures (images/):
%   fig_state_centroids_col_small.pdf, fig_state_pca_opt_col_small_v2.pdf   (states, half column each)
%   fig_balance_gem5_column_leela.pdf                                       (balance case, single column)
%   fig_stage_jain_opt_cpi_column.pdf                                       (balance across methods, single column)
% Numbers: results_small/state_clusters_v1.txt, robust_clusters_small.txt, results_opt_clusters/independent_clusters_small.txt,
%          per-benchmark state spread from results_small/state_clusters_v1.npz (labels_<bench>),
%          gem5_val4/settings_validation_pooled.txt (state agreement), results_rw (balance).

\subsection{Characterizing the Design Space}
\label{sec:char_states}
\label{sec:char_balance}

A CPI stack explains one design.
%
With a stack for every searched design, two stack-based analysis tools can be applied to the whole design space: bottleneck-state clustering~\cite{pesto} and bottleneck balance~\cite{archexplorer}.
%
We use them to ask two questions: which bottleneck states does the space contain, and do better designs have more balanced bottlenecks?

\begin{figure}[t]
\centering
\begin{subfigure}[t]{0.49\columnwidth}
  \centering
  \includegraphics[width=\linewidth]{images/fig_state_centroids_col_small.pdf}
  \caption{Stall composition}
  \label{fig:state_centroids}
\end{subfigure}\hfill
\begin{subfigure}[t]{0.49\columnwidth}
  \centering
  \includegraphics[width=\linewidth]{images/fig_state_pca_opt_col_small_v2.pdf}
  \caption{States across search methods}
  \label{fig:state_pca}
\end{subfigure}
\caption{Bottleneck states. (a) Centroid of each state: stall share of FP/BP/SP at Dispatch (solid) and Commit (hatched); each stage accounts for about half. (b) Greedy designs on the first two principal components, colored by state; circles: state centroids; black markers: centroids from clustering BO, GA, and random-search designs independently.}
\label{fig:states}
\end{figure}

\paragraph{Tool 1: bottleneck states}
Following Pesto~\cite{pesto}, we describe each design by the stall shares of FP, BP, and SP at Dispatch and Commit, the entry and exit of the out-of-order window, and group designs with k-means, choosing $k$ by the elbow method.
%
Pesto clusters the instructions of one simulated design; we cluster the 337k designs of greedy search across all benchmarks.

The elbow yields six states (Fig.~\ref{fig:state_centroids}).
%
They differ mainly at Dispatch, ordered from frontend- to backend-bound: BP accounts for 11\%, 23\%, 37\%, and 47\% of Dispatch stalls in S0, S2, S1, and S3.
%
S4 carries the most squash stalls, and S5 is rare (0.6\%).
%
This order is the dominant axis of the space: the first principal component explains 79\% of the variance (Fig.~\ref{fig:state_pca}).
%
Commit is nearly identical across states, with 89--99\% of its stalls being BP in all states but S5, a property of how Commit attributes stalls rather than of the designs.

The states are not an artifact of the search, the workload, or the model.
%
Clustering the designs of BO, GA, and random search independently yields five of the same states and no new one (Fig.~\ref{fig:state_pca}, black markers).
%
Each of S0--S4 occurs in 5 to 11 benchmarks, and 15 of 17 benchmarks span at least two states.
%
PRISM and gem5 assign the same or an adjacent state to 85\% of the validated designs (\S\ref{sec:char_settings}), and re-clustering resampled searches reproduces the partition (ARI 0.96--0.99).

\paragraph{Tool 2: bottleneck balance}
Bottleneck-driven DSE~\cite{archexplorer} repeatedly relieves the largest bottleneck, which implies that better designs have more balanced bottlenecks.
%
We test this premise at the granularity of stall categories, measuring the balance of stage $s$ with the Jain index~\cite{jain1984} of its stall cycles,
\begin{equation}
J_s = \frac{\left(\sum_{k} c_{k,s}\right)^2}{3\sum_{k} c_{k,s}^2}, \quad k \in \{\mathrm{FP}, \mathrm{BP}, \mathrm{SP}\},
\end{equation}
which ranges from $1/3$ (a single stall category) to 1 (equal categories).

\begin{figure}[t]
\centering
\includegraphics[width=\columnwidth]{images/fig_stage_jain_opt_cpi_column.pdf}
\caption{Per-stage balance (left axis) and CPI normalized to the search start (right axis, dotted) along four search methods; median over 17 benchmarks. For BO, GA, and random search, the design shown is the best one found so far within the area budget.}
\label{fig:balance_opt}
\end{figure}

\begin{figure}[t]
\centering
\includegraphics[width=\columnwidth]{images/fig_balance_gem5_column_leela.pdf}
\caption{Dispatch balance along greedy search on four benchmarks. Model and gem5 are evaluated on the same representative search; median and 25--75\% are over 1,000 searches.}
\label{fig:balance_case}
\end{figure}

Along the searches, CPI decreases by 13--37\%, yet balance barely moves (Fig.~\ref{fig:balance_opt}): for BO, GA, and random search, Dispatch balance changes by at most 0.04, and only greedy search shows a clear rise (0.58$\rightarrow$0.72).
%
Statically, among 100k random designs at matched area, the best 2\% have a Dispatch balance only 0.06 higher than the rest.
%
Dispatch is the most balanced stage under every method, while Commit stays near $1/3$, again because of its accounting.

gem5 confirms these trends on the same designs (Fig.~\ref{fig:balance_case}): balance rises on nab, saturates on leela, stays low on gcc, and on cactuBSSN peaks midway and falls back (model 0.68$\rightarrow$0.52, gem5 0.64$\rightarrow$0.48) while CPI keeps decreasing.
%
Balance is thus weakly associated with design quality but is not a monotone proxy for it: the premise of bottleneck-driven DSE does not hold at the category level.

\paragraph{Summary}
Across the space, bottlenecks fall into a few states ordered along one frontend--backend axis, and balance does not reveal whether further tuning pays off.
%
Next, we follow a search that stalls and ask whether it is trapped in a state and what pressure remains (\S V-C).
